\section{Introduction}

A bound on the matrix multiplication exponent and a finite description of
an algorithm are closely related, but they are different outputs of a
proof. An asymptotic tensor argument may establish the existence of a
suitable rank decomposition without giving its dimension or coefficients.
A finite extraction identifies the tensor maps, chooses their parameters,
and accounts for every cost needed to turn them into an arithmetic
algorithm. This paper carries out that extraction for the convolution
relations used in the preprint \emph{An Upper Bound of \(9/4\) for the
Matrix Multiplication Exponent}~\citep{OAI26}.

That preprint introduces a separation of tensor blocks that share one leg,
applies it to polynomial multiplication, and derives two inequalities for
a symmetrized convolution profile. The profile grows at least as
\(a^{4/3}\) on the diagonal. A detecting-character argument then gives
the claimed bound \(\omega\le9/4\), where \(\omega\) is the infimum
of arithmetic exponents for square matrix multiplication. The last
algorithmic step chooses a matrix size and an exact rank decomposition
whose existence follows from that bound. It does not specify the chosen
decomposition.

We replace this last passage by a direct construction. The starting point
is an exact solution of the scalar profile problem. Its affine rows
determine rational probabilities for the two tensor relations. Explicit
positive rational multipliers make these relations telescope, leaving a
diagonal term and a boundary term in the formal resource identity.
Clearing denominators turns this identity into integer tensor
types. We then construct the associated maps and derive a recursive
algorithm with all finite costs included.

\subsection{Main results}

The principal conclusion is the following effective form of the exponent
bound. A rational bilinear seed means a finite list of rational linear
forms whose products and recombinations compute a square matrix product.

\begin{theorem}\label{thm:main}
For every positive rational \(\eps\), a terminating procedure produces
an integer \(d\ge2\) and an exact rational bilinear algorithm for
\(d\times d\) matrix multiplication using fewer than
\(d^{9/4+\eps}\) bilinear multiplications, each multiplying a linear
form in the left input by a linear form in the right input.
A fixed recursion on this seed
multiplies \(n\times n\) matrices over any field of characteristic zero
in \(O_\eps(n^{9/4+\eps})\) arithmetic operations.
\end{theorem}

In particular, \(\omega\le9/4\). A bound with any positive real
allowance follows by choosing a smaller positive rational allowance.

The procedure first gives a relative algorithm: a fixed number of larger
matrix products are reduced simultaneously to a fixed number of smaller
ones. A padded batch recursion makes this a usual arithmetic algorithm.
Unrolling that recursion and restricting coefficients from its fixed
cyclotomic field to \(\Q\) gives the seed in Theorem~\ref{thm:main}.
The parameter choices use exact integer comparisons, so the statement
does not depend on a numerical tolerance or an unspecified rank search.
The complete proof is in Theorem~\ref{thm:effective} and
Corollary~\ref{cor:seed}.
For allowances at most \(3/4\), the prescribed seed generator returns a
dimension larger than the relative scheme's recursive input dimension.

Three intermediate results describe both the construction and its limits.
First, the least positive symmetric profile satisfying normalization,
discrete concavity, and shifted tripling is
\[
 F(m,n)=s_{\min(m,n)}\left(\max(m,n)+\frac{\min(m,n)-1}{2}\right),
 \qquad s_a=\prod_{h=1}^{a-1}\left(1+\frac1{3h}\right).
\]
This gives the exact scalar bound at every pair of indices, rather than
only its diagonal growth rate. Second, we give closed formulas for the
positive rational multipliers and prove the resulting integer resource
identity. Third, we determine their minimal common denominator and show
that it grows at least exponentially with the profile index.
The diagonal certificate uses only tripling at the row index; the full
least-profile theorem allows all shifted tripling inequalities.

For fixed index \(a\), the constructed exponents approach
\[
 \tau_a=\frac{3\log(2a-1)}{\log F(a,a)}.
\]
Every \(\tau_a\) is strictly larger than \(9/4\), and these values
tend to \(9/4\) as \(a\) increases. At each finite stage the cost
formula retains the source copies, output multiplicity, interpolation,
catalyst supply, and rounding of the recursive branch count. We also
reduce the shared-leg separator from the preprint's \(5M\) source
copies to \(2M-1\), and give a separate exact two-power projection.

The numerical exponent \(9/4\) is already the conclusion announced in
\citet{OAI26}. Our contribution is the explicit finite extraction and
its accompanying sharp scalar and finite-cost analysis. The proof below
is self-contained at the level of tensor maps. It uses neither the
detecting-character theorem nor the spectral appendix of that preprint.

\subsection{Relation to previous work}

Strassen's recursive use of a finite bilinear identity established the
basic link between small matrix algorithms and an arithmetic exponent
below three~\citep{Str69}. Approximate bilinear algorithms and their
conversion to exact algorithms by tensor powers and interpolation form
another part of the foundation~\citep{Bini80}.
Sch\"onhage's asymptotic sum inequality relates the simultaneous
computation of independent matrix products to the exponent~\citep{Sch81}.
An accessible detailed account of finite interpolation and the sum
inequality is given by \citet[Theorem~7 and its preceding
discussion]{Sto10}. Our root-of-unity averaging is an explicit instance
of this interpolation principle. Our batch recursion keeps the
simultaneous-output count visible instead of invoking the sum inequality
as a black box.

The asymptotic spectrum provides a general language for comparing tensor
powers through numerical invariants~\citep{Str88,CVZ}. Relative bilinear
complexity is likewise an established setting for reductions between
matrix multiplication tensors~\citep{Str87}. Catalytic tensor comparisons
also occur in the study of matrix multiplication: an auxiliary tensor is
retained on both sides of a reduction~\citep{CVZ21}. The present construction
starts from specific finite relations within this framework. Its rational
certificate supplies the multiplicities needed to compose those relations,
so the proof of the algorithm does not require an abstract separation
argument.

Fixed-frequency restrictions and entropy estimates are standard in
tensor-power methods. Le Gall's analysis of tensor powers gives a detailed
example of these methods and of balancing products over leg
permutations~\citep{LG14}. Recent upper bounds based on refinements of
the laser method also rely on carefully chosen distributions and
asymmetric tensor powers~\citep{AW21,AD25,Dup26}.
For numerical context, the SODA 2025 paper of \citet{AD25} proves
\(\omega<2.371339\), while the 2026 preprint of \citet{Dup26}
reports \(\omega<2.371177\). The \(9/4\) bound announced in the
October 2026 preprint \citep{OAI26} is substantially stronger; our proof
establishes the finite tensor route to that bound independently of the
preprint's spectral conclusion. The construction uses a specific family
of convolution relations and a closed rational certificate for it.
Interpolation, entropy counting, symmetrization, and recursive use of
bilinear identities are established tools.

The closest source for the actual tensor maps is \citet{OAI26},
Proposition~3.1 and Section~4. We give the bases and coordinate maps
explicitly, strengthen its two initial degenerations to exact type
restrictions, and improve its separator count. The least-profile formula,
the rational telescoping weights, their denominator formula, and the
complete finite cost expression are the specific developments established
here. These are comparisons with an identified prior argument; no claim
of a first general constructive conversion of asymptotic tensor bounds
is needed.

In particular, an existence theorem for a suitable complex rank
decomposition already implies an eventual symbolic search: its defining
polynomial equations have rational coefficients, and a complex solution
can be replaced by an algebraic one. The present procedure has a different
purpose. It generates identities by prescribed maps, with formulas for
their finite resource counts, rather than enumerating candidate solutions
to rank equations. Its value is mathematical control of that passage,
not a practical search bound.

\subsection{Scope and organization}

The construction concerns exact arithmetic. It has very large fixed
constants, no established numerical stability bound, and no demonstrated
useful matrix-size range. The cost analysis gives a reason for this size:
the exact rational types alone have exponentially growing normalization.
For the prescribed index-two batch schedule, even parameters that beat
\(\log_4 48\) require a recursive input dimension with more than
11.37 billion decimal digits. This is a limitation of the specified
schedule, not of arbitrary algorithms with the same exponent.

Section~\ref{sec:preliminaries} fixes the tensor conventions and the
arithmetic model. Section~\ref{sec:profile} proves the least-profile
theorem and the rational certificate. Section~\ref{sec:maps} gives the
type maps and separator. Section~\ref{sec:extraction} composes them,
proves the batched recursion, and specifies terminating parameter choices.
Section~\ref{sec:costs} analyzes the denominator growth and finite
example. Section~\ref{sec:discussion} describes the consequences
and the limits of the method. The appendices specify the full composed
relation maps, prove the optional separator refinements, and give the
detailed lower bound for the small-index schedule.
